\documentclass[aspectratio=169]{beamer}
\usetheme{Madrid}
\usepackage[spanish,es-nodecimaldot]{babel}
\usepackage[utf8]{inputenc}
\usepackage[T1]{fontenc}
\usepackage{amsmath}
\usepackage{amssymb}
\usepackage{graphicx}
\usepackage{booktabs}
\usepackage{xcolor}

% Color UNS
\definecolor{azulUNS}{RGB}{30,60,120}
\usecolortheme[named=azulUNS]{structure}

\title[MAD1 -- U6]{Métodos de Análisis de Datos I}
\subtitle{Unidad 6 -- Bloque 1: Otros paradigmas de aprendizaje}
\author{Departamento de Matemática -- UNS}
\date{2026}

\begin{document}

\begin{frame}
  \titlepage
\end{frame}

\begin{frame}{Contenidos}
  \tableofcontents
\end{frame}

%-----------------------------------------------
\section{El mapa del aprendizaje automático}
%-----------------------------------------------

\begin{frame}{¿Dónde estamos?}
  Durante el curso trabajamos principalmente con:
  \begin{itemize}
    \item \textbf{Aprendizaje supervisado:} regresión, clasificación, SVM (Support Vector Machine, máquina de vectores de soporte), árboles, redes neuronales
    \item \textbf{Aprendizaje no supervisado:} clustering, reducción de dimensionalidad, detección de anomalías, autoencoders
  \end{itemize}
  \bigskip
  Pero el ecosistema de ML (Machine Learning, aprendizaje automático) es mucho más amplio. En este bloque hacemos un mapa general de los paradigmas existentes.
\end{frame}

\begin{frame}{Los paradigmas: una visión de conjunto}
  \begin{columns}
    \column{0.5\textwidth}
      \textbf{Por el tipo de señal de entrenamiento:}
      \begin{itemize}
        \item Supervisado
        \item No supervisado
        \item Semi-supervisado
        \item Auto-supervisado
        \item Por refuerzo
      \end{itemize}
    \column{0.5\textwidth}
      \textbf{Por otras características:}
      \begin{itemize}
        \item Aprendizaje federado
        \item Aprendizaje activo
        \item Few-shot / zero-shot
        \item Transfer learning
      \end{itemize}
  \end{columns}
  \bigskip
  \begin{alertblock}{Idea clave}
    No son categorías mutuamente excluyentes. Un modelo puede ser a la vez auto-supervisado y usar transfer learning.
  \end{alertblock}
\end{frame}

%-----------------------------------------------
\section{Aprendizaje semi-supervisado}
%-----------------------------------------------

\begin{frame}{Aprendizaje semi-supervisado}
  \begin{block}{Definición}
    Se cuenta con un conjunto de datos \textbf{mixto}: pocos ejemplos etiquetados y muchos sin etiquetar. El modelo aprovecha la estructura de los datos no etiquetados para mejorar la clasificación.
  \end{block}
  \bigskip
  \textbf{¿Qué lo diferencia?}
  \begin{itemize}
    \item En supervisado: todos los datos tienen etiqueta.
    \item En no supervisado: ningún dato tiene etiqueta.
    \item En semi-supervisado: \textit{algunos} tienen etiqueta, y eso es valioso.
  \end{itemize}
\end{frame}

\begin{frame}{Semi-supervisado: ejemplos}
  \small
  \begin{itemize}
    \item \textbf{Imágenes médicas:} etiquetar una radiografía requiere un radiólogo (caro y lento). Con 200 imágenes etiquetadas y 10.000 sin etiquetar, el modelo usa la estructura de las no etiquetadas para ubicar mejor la frontera entre tejido sano y patológico.
    \item \textbf{Detección de spam:} clasificar correos a mano es tedioso. Unos pocos miles etiquetados más millones sin etiquetar permiten aprender qué patrones de lenguaje ``se parecen'' a spam, sin anotar todo el conjunto.
    \item \textbf{Reconocimiento de voz:} transcribir audio palabra por palabra es costoso. Miles de horas de audio sin transcribir ayudan a aprender la fonética del idioma; las pocas horas transcritas anclan el significado.
    \item \textbf{Análisis de sentimientos:} reseñas con puntaje (etiqueta implícita: 1 a 5 estrellas) más millones de comentarios sin calificación.
  \end{itemize}
  \begin{exampleblock}{Estrategia típica: \textit{self-training}}
    Entrenar con los etiquetados $\rightarrow$ predecir sobre los no etiquetados $\rightarrow$ agregar las predicciones más confiables al conjunto de entrenamiento $\rightarrow$ repetir.
  \end{exampleblock}
\end{frame}

%-----------------------------------------------
\section{Aprendizaje auto-supervisado}
%-----------------------------------------------

\begin{frame}{Aprendizaje auto-supervisado}
  \begin{block}{Definición}
    El modelo \textbf{genera sus propias etiquetas} a partir de los datos crudos. No necesita anotación humana. La tarea de entrenamiento es artificial (\textit{pretext task}), pero lo aprendido es útil para tareas reales.
  \end{block}
  \bigskip
  \textbf{¿Qué lo diferencia del no supervisado?}
  \begin{itemize}
    \item No supervisado: descubre estructura (clusters, dimensiones).
    \item Auto-supervisado: se entrena \textit{como si} fuera supervisado, pero sin anotación humana.
  \end{itemize}
\end{frame}

\begin{frame}{Auto-supervisado: ejemplos (1/2) --- la tarea pretexto}
  La clave es inventar una \textbf{tarea pretexto}: un problema cuya respuesta ya está en los datos, sin que nadie la anote. Ejemplos en imágenes:
  \bigskip
  \begin{itemize}
    \item \textbf{Colorización:} se pasa la imagen a blanco y negro y se entrena al modelo a recuperar el color. La etiqueta (el color real) estaba en la foto original.
    \item \textbf{Inpainting (relleno):} se oculta un parche de la imagen y el modelo debe predecir qué había ahí. Para hacerlo bien, debe ``entender'' el contenido.
    \item \textbf{Predicción de rotación:} se rota la imagen 0\textdegree, 90\textdegree, 180\textdegree{} o 270\textdegree{} y el modelo debe adivinar el ángulo. Reconocer la orientación correcta exige aprender qué es cada objeto.
  \end{itemize}
  \bigskip
  En todos los casos: \textbf{no hubo anotación humana}, pero el modelo aprende representaciones útiles para tareas reales posteriores.
\end{frame}

\begin{frame}{Auto-supervisado: ejemplos (2/2) --- modelos reales}
  \small
  \begin{itemize}
    \item \textbf{Modelos de lenguaje --- GPT (Generative Pre-trained Transformer) y BERT (Bidirectional Encoder Representations from Transformers):} la tarea pretexto es predecir la próxima palabra (GPT) o rellenar una palabra oculta en la oración (BERT). La etiqueta es el propio texto. Ambos son \textbf{arquitecturas} basadas en Transformers.
    \item \textbf{Visión por computadora --- SimCLR y DINO:} son \textbf{métodos} de entrenamiento contrastivo. Se generan dos versiones distorsionadas de la misma imagen (recorte, color, giro) y se entrena al modelo para que las reconozca como ``la misma cosa''.
    \item \textbf{Audio --- wav2vec:} aprende a predecir fragmentos de audio a partir de los anteriores, sin transcripciones. Luego ese conocimiento se reutiliza para reconocimiento de voz.
  \end{itemize}
  \begin{alertblock}{Relevancia}
    La mayoría de los grandes modelos de IA actuales (LLMs --- Large Language Models, grandes modelos de lenguaje; y modelos de visión fundacionales) son auto-supervisados. Es el paradigma dominante en deep learning moderno.
  \end{alertblock}
\end{frame}

%-----------------------------------------------
\section{Aprendizaje por refuerzo}
%-----------------------------------------------

\begin{frame}{Aprendizaje por refuerzo (RL)}
  \begin{block}{Definición}
    Un \textbf{agente} aprende a tomar decisiones en un \textbf{entorno} mediante prueba y error, maximizando una \textbf{recompensa} acumulada. No hay un conjunto de datos fijo: los datos se generan por la interacción.
  \end{block}
  \bigskip
  \textbf{Componentes clave:}
  \begin{itemize}
    \item \textbf{Estado} $s$: descripción del entorno en un momento dado.
    \item \textbf{Acción} $a$: lo que el agente puede hacer.
    \item \textbf{Recompensa} $r$: señal escalar que indica qué tan buena fue la acción.
    \item \textbf{Política} $\pi$: función que mapea estados a acciones.
  \end{itemize}
\end{frame}

\begin{frame}{RL: ¿qué lo diferencia de todo lo demás?}
  \begin{columns}
    \column{0.5\textwidth}
      \textbf{En supervisado/no supervisado:}
      \begin{itemize}
        \item Datos fijos, estáticos.
        \item No hay interacción con un entorno.
        \item El error es directo (pérdida, distancia).
      \end{itemize}
    \column{0.5\textwidth}
      \textbf{En RL:}
      \begin{itemize}
        \item Los datos se generan al actuar.
        \item Las acciones afectan el entorno (y futuros estados).
        \item La recompensa puede ser diferida.
      \end{itemize}
  \end{columns}
  \bigskip
  \begin{exampleblock}{La gran dificultad}
    El \textbf{problema de atribución de crédito}: ¿qué acción, entre muchas pasadas, fue responsable de una recompensa que llegó tarde?
  \end{exampleblock}
\end{frame}

\begin{frame}{RL: ejemplos (1/2)}
  \small
  \begin{itemize}
    \item \textbf{Juegos:} AlphaGo (Go), AlphaStar (StarCraft II), OpenAI Five (Dota 2). Son entornos ideales para RL: reglas claras y una recompensa inequívoca (ganar o perder). El agente juega millones de partidas contra sí mismo y mejora por ensayo y error.
    \item \textbf{Robótica:} aprender a caminar, agarrar objetos o ensamblar piezas. La recompensa premia el movimiento exitoso; el agente entrena en simulación porque equivocarse en el mundo real rompe hardware.
    \item \textbf{Control de tráfico:} semáforos que ajustan sus tiempos de verde según el flujo de autos que detectan, con la recompensa definida como minimizar la demora total. El entorno (la ciudad) responde a cada decisión.
  \end{itemize}
\end{frame}

\begin{frame}{RL: ejemplos (2/2)}
  \small
  \begin{itemize}
    \item \textbf{Control industrial:} optimizar el consumo energético de la refrigeración en centros de datos (Google DeepMind logró reducir el gasto de enfriamiento). La recompensa es la energía ahorrada sin violar límites de temperatura.
    \item \textbf{Sistemas de recomendación:} qué contenido mostrar para maximizar el interés \textit{a largo plazo}, no solo el clic inmediato. El reto es que la acción de hoy cambia qué verá el usuario mañana.
    \item \textbf{Finanzas:} estrategias de trading y gestión de cartera, donde cada operación modifica el estado del portafolio.
    \item \textbf{RLHF (Reinforcement Learning from Human Feedback):} los LLM como ChatGPT se afinan con RL usando preferencias humanas como señal de recompensa, para alinear las respuestas con lo que las personas consideran útil.
  \end{itemize}
\end{frame}

%-----------------------------------------------
\section{Otros paradigmas relevantes}
%-----------------------------------------------

\begin{frame}{Aprendizaje federado}
  \begin{block}{Definición}
    El modelo se entrena \textbf{de forma distribuida}: los datos nunca salen de cada dispositivo/nodo. Solo se comparten los gradientes o actualizaciones del modelo.
  \end{block}
  \bigskip
  \begin{itemize}
    \item \textbf{Motivación:} privacidad, regulación (datos médicos, financieros), ancho de banda.
    \item \textbf{Ejemplo:} el teclado predictivo de tu celular mejora con el uso de millones de usuarios sin que Google/Apple vea tus mensajes.
    \item \textbf{Desafíos:} datos heterogéneos entre nodos, comunicación costosa, ataques adversariales.
  \end{itemize}
\end{frame}

\begin{frame}{Aprendizaje activo}
  \begin{block}{Definición}
    El modelo \textbf{elige qué ejemplos quiere que le etiqueten}. Consulta a un oráculo (humano) sobre los casos más informativos o inciertos.
  \end{block}
  \bigskip
  \begin{itemize}
    \item \textbf{Motivación:} etiquetar datos es caro. Con menos etiquetas bien elegidas se logra el mismo rendimiento.
    \item \textbf{Estrategias:} mayor incertidumbre, mayor divergencia entre versiones del modelo (\textit{query by committee}), mayor densidad informativa.
    \item \textbf{Ejemplo:} clasificación de imágenes satelitales de desastres — el modelo pide etiquetar primero los casos ambiguos.
  \end{itemize}
\end{frame}

\begin{frame}{Few-shot y zero-shot learning}
  \begin{block}{Few-shot}
    El modelo aprende una nueva tarea con \textbf{muy pocos ejemplos} (1, 5, 10). Requiere haber aprendido representaciones generales durante el preentrenamiento.
  \end{block}
  \begin{block}{Zero-shot}
    El modelo resuelve una tarea \textbf{sin ver ningún ejemplo} de esa tarea en entrenamiento. Solo se describe la tarea en lenguaje natural.
  \end{block}
  \bigskip
  \begin{exampleblock}{Ejemplo}
    ``Clasificá este texto como positivo o negativo.'' --- un LLM puede hacer esto sin haber sido entrenado explícitamente para análisis de sentimientos.
  \end{exampleblock}
\end{frame}

\begin{frame}{Transfer learning}
  \begin{block}{Definición}
    Se toma un modelo preentrenado en una tarea grande y se \textbf{adapta} a una tarea más pequeña o específica (\textit{fine-tuning}).
  \end{block}
  \bigskip
  \begin{itemize}
    \item \textbf{Ejemplo clásico:} ResNet (Residual Network, una arquitectura de red convolucional) entrenada en ImageNet (un dataset de millones de imágenes) $\rightarrow$ clasificar rayos X con apenas 1.000 imágenes.
    \item \textbf{Ejemplo en NLP (Natural Language Processing, procesamiento de lenguaje natural):} BERT preentrenado en Wikipedia $\rightarrow$ clasificar reseñas médicas.
    \item \textbf{Por qué funciona:} las capas bajas aprenden representaciones generales (bordes, texturas, sintaxis) que son útiles en muchos dominios.
  \end{itemize}
  \bigskip
  \textbf{Transfer learning} + \textbf{auto-supervisado} = la base de la IA moderna.
\end{frame}

%-----------------------------------------------
\section{Diagrama de Venn: relaciones entre paradigmas}
%-----------------------------------------------

\begin{frame}{¿Cómo se relacionan?}
  \begin{columns}
    \column{0.55\textwidth}
      \textbf{Solapamientos frecuentes:}
      \begin{itemize}
        \item Auto-supervisado $\cap$ Transfer learning: GPT, BERT, CLIP (Contrastive Language-Image Pre-training, modelo multimodal que liga texto e imagen).
        \item Semi-supervisado $\cap$ Aprendizaje activo: etiquetar selectivamente + propagar.
        \item RL $\cap$ Auto-supervisado: RLHF en LLMs.
        \item Federado $\cap$ cualquier paradigma: es una restricción sobre \textit{dónde} viven los datos, no sobre \textit{cómo} se aprende.
      \end{itemize}
    \column{0.45\textwidth}
      \textbf{Lo que los diferencia en esencia:}
      \begin{itemize}
        \item \textbf{Supervisado:} etiquetas humanas.
        \item \textbf{No supervisado:} sin etiquetas.
        \item \textbf{Semi:} pocas etiquetas.
        \item \textbf{Auto:} etiquetas sintéticas.
        \item \textbf{RL:} recompensa por acción.
        \item \textbf{Activo:} el modelo pregunta.
        \item \textbf{Federado:} datos distribuidos.
      \end{itemize}
  \end{columns}
\end{frame}

%-----------------------------------------------
\section{Glosario de siglas y modelos}
%-----------------------------------------------

\begin{frame}{Glosario (1/2): conceptos y siglas generales}
  \small
  \begin{table}
    \begin{tabular}{lll}
      \toprule
      \textbf{Sigla} & \textbf{Significado} & \textbf{Qué es} \\
      \midrule
      ML   & Machine Learning                       & Disciplina \\
      RL   & Reinforcement Learning                 & Paradigma de aprendizaje \\
      NLP  & Natural Language Processing            & Disciplina \\
      LLM  & Large Language Model                   & Categoría de modelo \\
      SVM  & Support Vector Machine                 & Algoritmo \\
      RLHF & Reinf. Learning from Human Feedback    & Técnica de entrenamiento \\
      \bottomrule
    \end{tabular}
  \end{table}
  \bigskip
  Estas siglas nombran \textit{ideas} o \textit{enfoques}, no modelos concretos.
\end{frame}

\begin{frame}{Glosario (2/2): modelos, arquitecturas y sistemas}
  \small
  \begin{table}
    \begin{tabular}{lll}
      \toprule
      \textbf{Nombre} & \textbf{Significado} & \textbf{Qué es} \\
      \midrule
      GPT      & Generative Pre-trained Transformer       & Arquitectura (Transformer) \\
      BERT     & Bidir. Encoder Repr. from Transformers   & Arquitectura (Transformer) \\
      ResNet   & Residual Network                         & Arquitectura (convolucional) \\
      CLIP     & Contrastive Language-Image Pre-training  & Modelo multimodal \\
      SimCLR   & Simple Contrastive Learning of Repr.     & Método de entrenamiento \\
      DINO     & self-DIstillation with NO labels         & Método de entrenamiento \\
      wav2vec  & ``waveform to vector''                   & Modelo de audio \\
      AlphaGo  & ---                                      & Sistema/agente (RL) \\
      ImageNet & ---                                      & Dataset / benchmark \\
      \bottomrule
    \end{tabular}
  \end{table}
  \bigskip
  \begin{alertblock}{Tres ejes distintos}
    \textbf{Arquitectura} = cómo está construida la red (GPT, BERT, ResNet). \textbf{Método} = cómo se entrena (SimCLR, DINO, RLHF). \textbf{Sistema} = una aplicación completa (AlphaGo). Un proyecto real combina los tres.
  \end{alertblock}
\end{frame}



\begin{frame}{Ideas para llevarse}
  \begin{enumerate}
    \item Los paradigmas \textbf{no son cajas}: se combinan y se superponen.
    \item La escasez de etiquetas es el problema central que motiva semi-supervisado, activo y auto-supervisado.
    \item El aprendizaje por refuerzo es el más diferente: interacción, secuencia, recompensa diferida.
    \item El aprendizaje federado responde a restricciones del mundo real (privacidad, regulación), no a una diferencia en cómo se aprende.
    \item La IA de vanguardia hoy combina \textbf{auto-supervisado + transfer learning + RL}.
  \end{enumerate}
\end{frame}

\end{document}
